\section{Introduction}
\label{sec:intro}

The dynamic behavior of such circuits is naturally captured by
\emph{ordinary differential equations}~(ODEs).
Simpler RC and RL sub-circuits obey \emph{first-order} ODEs, while the
complete RLC system requires a \emph{second-order} ODE to describe the
interplay between energy stored in the inductor's magnetic field and the
capacitor's electric field.
Understanding the transition from first- to second-order models clarifies
why the resistor plays an indispensable role: it is the only purely
dissipative element in the circuit.

Our scientific question is:
\begin{quote}
  \itshape
  How do resistors influence the damping and stability of series RLC circuits
  as described by ODEs, and what are the practical implications in engineering
  applications?
\end{quote}

This paper answers that question through a two-part approach: first by
deriving ideal analytic solutions for RC, RL, and RLC circuits, and second by
comparing those ideal responses with numerical simulations that include
phenomenological non-ideal resistor and inductor behavior.

The remainder of this paper is structured as follows.
Section~\ref{sec:circuit} introduces the circuit configuration and the
governing physical principles.
Section~\ref{sec:methods} details the derivation of the governing equations,
identification of variables and parameters, and the analytic solutions for
each circuit type.
Section~\ref{sec:results} presents results from varying circuit parameters
and initial conditions.
Section~\ref{sec:nonideal} extends the ideal model to non-ideal components
and resonant forcing using numerical simulation.
Section~\ref{sec:discussion} interprets the findings in the context of our
research question, discusses limitations, and proposes future work.
Section~\ref{sec:conclusion} concludes the paper.
